\section{Decision three: what becomes of the impeded day}

\rg{93} names four outcomes. Which one applies is decided by rules of its own,
and it is a mistake --- the commonest one in this part of the code --- to infer
the outcome from the margin of defeat. A first-class feast beaten by a hair is
transferred; a second-class Patron beaten by a hair is commemorated or lost.

\subsection{Transfer: only first-class feasts, and only accidentally}

\begin{rulebox}{\rg{95} (first paragraph)}
Ius translationis in alium diem ob occurrentiam accidentalem cum die
liturgico, qui in tabella praecedentiae superiorem obtinet locum, competit
solummodo festis I classis. Alia festa, ab Officio gradus superioris
accidentaliter impedita, aut commemorantur aut, eo anno, penitus omittuntur,
iuxta rubricas.
\end{rulebox}

Where a first-class feast is transferred, it goes ``in proximum sequentem diem
qui non sit I vel II classis'' (\rg{96}) --- the next following day that is
neither first nor second class, which may be several days away. Two feasts are
excepted and given a fixed destination described as their own seat:

\begin{rulebox}{\rg{96}\,a--b}
a) festum Annuntiationis B. Mariae Virg., quando est transferendum post
Pascha, transfertur, tamquam in sedem propriam, in feriam II post dominicam in
albis; b) Commemoratio omnium Fidelium defunctorum, quando occurrit cum
dominica, transfertur, tamquam in sedem propriam, in feriam II sequentem.
\end{rulebox}

\noindent
\latin{Tamquam in sedem propriam} matters. A day reached as a proper seat is
not a borrowed lodging: on it the transferred item behaves as though it were at
home, and the third-class feast that ordinarily occupies that date is impeded
perpetually as far as that year is concerned. Both worked cases below turn on
this phrase.

Two further ordering rules cover pile-ups. If several first-class feasts fall
on one day, the highest in the table is kept and the others are transferred in
the order in which the table lists them (\rg{97}); if several must be
transferred onto following days, the same table order governs, and on a tie the
one impeded first goes first (\rg{98}). And \rg{99} settles a question that
would otherwise be open: ``Festa translata sunt eiusdem gradus ac in sede
propria'' --- a transferred feast keeps the class it had, with the privileges
that class carries.

\subsection{Reposition: perpetual collisions only}

\rg{100} grants the right of reposition ``ob occurrentiam perpetuam'' to all
first- and second-class feasts, and also to particular third-class feasts
falling outside Advent and Lent which are impeded throughout a whole diocese or
Order or in their own church. Everything else perpetually impeded is
perpetually commemorated or perpetually omitted. First- and second-class
feasts are repositioned to the nearest following day that is not first or
second class; third-class ones to the nearest following day free of Offices of
equal or higher grade (\rg{101}). And \rg{102} makes the new date genuinely
proper: ``Dies in quem festa perpetuo impedita reposita sunt, habetur tamquam
dies proprius.''

A neat trap sits in \rg{94}: ``Commemoratio die fixo statuta non transfertur
vel reponitur cum festo transferendo vel reponendo, sed fit suo die vel
omittitur.'' A commemoration attached to a date does not travel with the feast
it accompanies. It is made on its own date or lost.

\subsection{Commemoration, and what 1960 took away}

The commemoration apparatus is the part of the older system that the 1960 code
most drastically reduced, and it did so not by abolishing commemorations but by
capping their number and excluding whole categories.

\subsubsection*{Two kinds}

\begin{rulebox}{\rg{107}--\rg{108}}
Commemorationes sunt aut \textit{privilegiatae} aut \textit{ordinariae}.
Commemorationes \textit{privilegiatae} fiunt in Laudibus et in Vesperis necnon
in omnibus Missis; commemorationes vero \textit{ordinariae} fiunt tantum in
Laudibus, in Missis conventualibus et in omnibus Missis lectis.
\end{rulebox}

\noindent
Read \rg{108} twice. An ordinary commemoration is made at every Low Mass and
at the conventual Mass --- and at no other sung Mass. A sung parish Mass that
is not the conventual Mass of a choir obligation drops every ordinary
commemoration, whatever the day. \rg{111}\,a and \rgm{434}\,a say it again
from the other direction, and Section~6 returns to it, because it is the single
most-missed rule in the whole apparatus.

The privileged commemorations are a closed list of six:

\begin{rulebox}{\rg{109}}
Commemorationes \textit{privilegiatae} sunt commemorationes: a) de dominica;
b) de die liturgico I classis; c) de diebus infra octavam Nativitatis Domini;
d) de feriis Quatuor Temporum mensis septembris; e) de feriis Adventus,
Quadragesimae et Passionis; f) de Litaniis maioribus, in Missa. Omnes aliae
commemorationes sunt commemorationes \textit{ordinariae}.
\end{rulebox}

\noindent
Everything else --- every saint, every vigil except as it is a first-class day,
every second-class feast, every octave day of Christmas taken as a feast --- is
ordinary.

\subsubsection*{How many are admitted}

\begin{rulebox}{\rg{111}}
Ratio admittendi commemorationes haec est: a) in diebus liturgicis I classis
et in Missis in cantu non conventualibus, nulla admittitur commemoratio,
praeter unam privilegiatam; b) in dominicis II classis, una tantum admittitur
commemoratio, scilicet de festo II classis, quae tamen omittitur si
commemoratio privilegiata facienda sit; c) in aliis diebus liturgicis II
classis, una tantum admittitur commemoratio, scilicet aut una privilegiata aut
una ordinaria; d) in diebus liturgicis III et IV classis, duae tantum
admittuntur commemorationes.
\end{rulebox}

So the ceiling is: one privileged only on a first-class day; on a second-class
Sunday one commemoration and it must be of a second-class feast, displaced by
any privileged commemoration that is due; on other second-class days exactly
one of either kind; on third- and fourth-class days two. \rg{114} disposes of
the surplus without ceremony: ``Quaelibet commemoratio, quae numerum pro
singulis diebus liturgicis statutum superet, omittitur.''

\subsubsection*{Exclusions by identity, not by number}

\rg{112} then removes commemorations that would otherwise fit:

\begin{tablerule}{Excluded}{Rule}{Locus}
Same divine Person & An Office, Mass or commemoration of a feast or mystery of
one divine Person excludes a commemoration or prayer of another feast or
mystery of the same Person. & \rg{112}\,a\\
Sunday and feast of the Lord & An Office, Mass or commemoration of the
Sunday excludes a commemoration or prayer of a feast or mystery of the Lord,
and conversely. & \rg{112}\,b\\
Two of the season & An Office, Mass or commemoration \latin{de Tempore}
excludes another commemoration \latin{de Tempore}. & \rg{112}\,c\\
Duplicated intercession & An Office, Mass or commemoration of Our Lady or of a
saint or blessed excludes another commemoration or prayer in which the same
Lady's or saint's intercession is asked --- but this does not apply to the
prayer of the Sunday or feria in which the same saint is invoked. &
\rg{112}\,d\\
\end{tablerule}

One commemoration cannot be excluded at all. \rg{110} makes the mutual
commemoration of Ss.\ Peter and Paul \latin{inseparabilis}: in the Office or
Mass of either, the other's prayer is added under a single conclusion, and
``duae orationes adeo in unam coalescere censentur ut, in numero orationum
computando, pro unica habeantur.'' Two prayers counting as one.

\subsubsection*{Order}

\rg{113}: ``Commemoratio de Tempore fit primo loco. In admittendis et
ordinandis aliis commemorationibus, servetur ordo tabellae praecedentiae.''
The season first, then the table order --- and the Peter-and-Paul prayer, when
added \latin{ad modum commemorationis}, immediately after its partner and
before all others (\rg{110}\,c).

\begin{caution}{What the code removed outright}
The motu proprio's own third article revoked ``statuta, privilegia, indulta et
consuetudines cuiuscumque generis, etiam saecularia et immemorabilia \dots{}
quae his rubricis obstant.'' Concretely, on the evidence of the code and the
\work{Variationes} promulgated with it: octaves other than Christmas, Easter
and Pentecost cease to exist anywhere (\rg{64}); the \latin{Suffragium de
omnibus Sanctis} and the \latin{commemoratio de Cruce} are abolished, and the
seasonal second and third orations of Advent, Lent and the Sundays --- the
\latin{orationes pro diversitate Temporum} that older hand missals print
beneath the collect --- are abolished with them; the number of orations at Mass
may never exceed three ``nullo praetextu'' (\rgm{435}); and the imposed prayer
of the Ordinary is barred from every first- and second-class day and from every
sung Mass (\rgm{457}\,d). A reader holding a pre-1955 hand missal is therefore
holding a book whose commemoration apparatus has been dismantled, not merely
trimmed, and its printed second and third collects must not be used with the
1962 calendar.
\end{caution}
